\CharacterSheet{traitor}{Traidora}{

\shortdescription{Uma reformista compassiva, que prega os bons costumes.}

\ObjectivesHeading{Objetivos}
\begin{itemize}
  \item Descobrir qual era o veneno para matar os vampiros. (\PositiveReward{1}, \NegativeReward{1})
\end{itemize}

\LosingConditionsHeading{Condições de derrota}
Não conseguir expulsar por votação quem matou \Character{newcomer}, caso haja uma pessoa responsável pelo homicídio.

\SecretsHeading{Segredos}
\begin{itemize}
  \item Tenho trocado correspondência secreta com a \outsideentity{igreja}. Quero que os vampiros sejam curados, se possível, e destruídos, se não houver cura.
  \item Ontem, procurei no gabinete de \Character{usurper} uma forma de matar vampiros. Em vez disso, roubei e usei um ritual, espalhando objetos do gabinete pelo castelo.
\end{itemize}

\ActionsHeading{Ações}
\begin{itemize}
  \item Fala com carinho da dignidade humana, da entreajuda e da misericórdia. Oferece ajuda prática em vez de simpatia abstrata, sobretudo quando a festa precisar de voluntários.
\end{itemize}

\BackgroundHeading{Contexto}

Antes de ser vampira, passei a vida a ajudar pessoas ignoradas por quem tinha riqueza e poder. Organizava comida, abrigo e medicamentos onde conseguia. Acreditava que uma vida boa se media pela segurança que se criava para quem menos a tinha. Depois conheci \Character{lover}. Era forte sem ser cruel e fez-me acreditar que até alguém nascido para a violência podia escolher usá-la bem.

\timestamp{No ano 1400 d.C.}, \Character{lover} pediu a \Character{vampirehead} que me deixasse entrar na família. A decisão foi controversa. \Character{supremacist} opôs-se abertamente, e \Character{enforcer} não confiava numa humana que não parava de falar de misericórdia. \Character{vampirehead} aceitou-me porque \Character{lover} não desistia. Pensei que a imortalidade me daria séculos para fazer o bem. Em vez disso, deu-me séculos a ver esta família alimentar-se de humanos e chamar-lhe tradição.

Ainda me lembro daquela festa horrível, em que \Character{vampirehead}, \Character{usurper}, \Character{enforcer}, \Character{buttler}, \Character{ghoul}, \Character{supremacist}, \Character{lover} e \Character{prevlover} falavam das suas vidas. Como se riam do sofrimento dos outros.

Durante muito tempo, tentei mudar a família por dentro. Defendia a moderação, protegia quem podia e implorava aos outros que vissem os humanos como mais do que comida ou criados. A maioria ria-se de mim. Acabei por começar a trocar correspondência com a \outsideentity{igreja}. Procuram uma cura para o vampirismo. Se não houver cura, acredito que os vampiros que atacam humanos têm de ser travados de vez. Sei o que isso faz de mim aos olhos desta família. Prefiro que me chamem traidora a continuar a fingir que a crueldade é lealdade.

Uma vez, bêbeda, \Character{usurper} brincou com o facto de um vampiro morrer se receber um certo tratamento, enquanto me dizia para me matar. Raios. O que era? Não me lembro.

\timestamp{Há dois anos}, \Character{vampirehead} apresentou \Character{newcomer} num encontro de família. Percebi que a chefe o queria transformar e avisei-o discretamente de que a imortalidade nesta casa tinha um preço. Ouviu-me com educação, mas o poder e a sensação de pertença tinham o seu encanto.

Ao longo do último ano, tenho ficado inquieta com \Character{lover}. Reparei que guarda mensagens e desaparece a horas estranhas. Não sei se descobriu o meu contacto com a \outsideentity{igreja} ou se esconde algo seu. Ainda o amo, o que faz de cada incerteza uma ameaça.

\timestamp{Há dois meses}, \Character{vampirehead} anunciou que \Character{newcomer} se juntaria à família numa festa em \timestamp{16/7/2026}. Todos os vampiros tinham de comparecer. Umas duas semanas antes da festa, \Character{usurper} ofereceu-se para ajudar a preparar a celebração. \Character{vampirehead} perguntou então se mais alguém, além de \Character{usurper}, de \Character{buttler} e de \Character{newcomer}, queria chegar duas noites antes para ajudar. Recusei. Como toda a gente, percebi que os preparativos criariam uma oportunidade de entrar no gabinete de \Character{usurper}.

\timestamp{Ontem, à 01:00}, entrei nesse gabinete à procura da poção capaz de matar um vampiro. Vasculhei os frascos e as notas cifradas, mas não a encontrei. Em vez disso, levei um ritual chamado \textit{Fazer com que o que é deles seja de outra pessoa}. Acreditava que me permitiria obter a poção mais tarde. Não fiquei tempo suficiente para identificar os outros intrusos no gabinete.

\timestamp{Ontem, às 17:00}, depois de um estranho nevoeiro de memórias em falta se ter instalado na família, realizei o ritual. Já não conseguia identificar o frasco que queria, nem sequer enunciar claramente o meu objetivo. Parecia que não ia funcionar, até desaparecer de vista. Deveria ter deixado resíduos, mas não, nada.

\timestamp{Às 21:00}, percebi que \Character{lover} não estava em casa. Não o segui.

\GMNotes{
\paragraph{Porque se pode suspeitar desta personagem?}
\Character{traitor} trabalha em segredo com a \outsideentity{igreja}, condena abertamente a família, procurou meios de matar vampiros e realizou um ritual que espalhou provas do gabinete de \Character{usurper}. O seu motivo é real e pode ser exposto pela correspondência intercetada. É, por isso, uma falsa pista forte e deliberada.

\paragraph{Como se pode provar a inocência desta personagem?}
\Character{traitor} não encontrou a poção proibida, não mexeu no vinho ao almoço e realizou o ritual às \timestamp{17:00}, cinco horas depois de \Character{newcomer} já ter recebido a primeira dose escondida. O seu ritual explica por que razão os objetos do gabinete e as pistas estão espalhados pelo castelo, mas não administrou a dose que tornou fatal a transformação oficial.

\paragraph{Qual é o papel desta personagem na história e no jogo?}
\Character{traitor} oferece uma suspeita alternativa com profundidade moral. A jogadora pode obrigar o grupo a confrontar o tratamento que a família dá aos humanos, enquanto decide quanto revela de uma conspiração real. O ritual falhado também explica, dentro da história, por que estão as pistas espalhadas fora do gabinete de \Character{usurper}.

\paragraph{Como pode esta personagem encaminhar a investigação para a solução?}
Quando o enigma dos resíduos do ritual de roubo for resolvido, \Character{traitor} pode admitir que as poções e os papéis espalhados vieram do gabinete de \Character{usurper}. Isso incentiva o grupo a seguir as garrafas importantes até à sua dona. Combinado com o relato de \Character{buttler} de que \Character{usurper} substituiu o vinho ao almoço, transforma as aparentes provas do crime de \Character{traitor} num caminho para descobrir o plano homicida de \Character{usurper}.
}
}
